% Generated draft: requires mathematical and source review.
\documentclass[11pt]{article}
\usepackage[margin=25mm]{geometry}
\usepackage{fontspec}
\setmainfont{Latin Modern Roman}
\usepackage{amsmath,amssymb,mathtools,unicode-math}
\setmathfont{Latin Modern Math}
\usepackage{enumitem,needspace,xurl,hyperref}
\setlength{\parindent}{0pt}
\setlength{\parskip}{5pt}
\setlength{\emergencystretch}{3em}
\newcommand{\N}{\mathbb{N}}
\newcommand{\Z}{\mathbb{Z}}
\newcommand{\Q}{\mathbb{Q}}
\newcommand{\R}{\mathbb{R}}
\newcommand{\C}{\mathbb{C}}
\newcommand{\F}{\mathbb{F}}
\newcommand{\A}{\mathbb{A}}
\newcommand{\PP}{\mathbb{P}}
\newcommand{\E}{\mathbb{E}}
\newcommand{\eps}{\varepsilon}
\newcommand{\dd}{\,\mathrm{d}}
\begin{document}
\section*{UnsolvedMath ID: 30002098}
\textbf{Source problem code:} catalogue entry
\subsection{Definitions and mathematical statement}
The usual interpretation of a level-$\ell$ curve $[C,\eta]$ is a smooth, connected projective curve $C$ together with a nontrivial line bundle $\eta\in\operatorname{Pic}^{0}(C)$ of exact finite order $\ell\geq2$: thus $\eta^{\otimes\ell}\cong\mathcal O_C$ and $\eta^{\otimes m}\not\cong\mathcal O_C$ for $1\leq m<\ell$. The supplied record does not specify the ground field, its characteristic, the permitted levels, or the precise moduli convention; this interpretation therefore requires confirmation.

Let $g$ be the genus of $C$, let $K_C$ denote its canonical line bundle, and put $L=K_C\otimes\eta$ and $V=H^0(C,L)$. Under the usual interpretation, $\dim V=g-1$. When the complete linear series of $L$ defines an embedding, its paracanonical embedding is
\[
\varphi_L:C\hookrightarrow\mathbb P(V^{\vee})\cong\mathbb P^{g-2}.
\]
Let $S=\operatorname{Sym}V$ with its standard grading, and let $A=S/I_{\varphi_L(C)}$ be the homogeneous coordinate ring. Its minimal graded free resolution is called pure if each nonzero free module in homological degree $i$ has the form $S(-d_i)^{\beta_i}$ for a single degree $d_i$.

For $p\geq1$, define $K_{p,1}(C,L)$ as the middle cohomology of the Koszul complex
\[
\bigwedge^{p+1}V\otimes H^0(C,\mathcal O_C)
\xrightarrow{d}
\bigwedge^pV\otimes H^0(C,L)
\xrightarrow{d}
\bigwedge^{p-1}V\otimes H^0(C,L^{\otimes2}),
\]
where $d$ is the alternating differential induced by multiplication of sections.

The supplied question is: for a general level-$\ell$ curve $[C,\eta]$ of every even genus $g\geq6$, does the paracanonically embedded curve have a pure minimal graded resolution of its homogeneous coordinate ring? The record presents this as equivalent to
\[
K_{g/2-2,1}(C,K_C\otimes\eta)=0.
\]
Here ``general'' refers to a dense Zariski-open subset of the intended moduli space of level curves. That moduli space, the admissible levels and ground fields, and any exceptional cases are not specified in the record. The embedding and the asserted equivalence are retained as source-level qualifications; their hypotheses require mathematical review.
\subsection{Short English statement}
For a general curve of even genus $g\geq6$ equipped with a torsion line bundle $\eta$ of level $\ell$, ask whether its paracanonical homogeneous coordinate ring has a pure resolution. The supplied record states the equivalent target $K_{g/2-2,1}(C,K_C\otimes\eta)=0$. The ground field, level restrictions, and possible exceptions remain unresolved.
\subsection{Sources}
\begin{itemize}
\item Local registry: \texttt{data/unsolvedmath/problems.json}, permanent UnsolvedMath ID $30002098$, source problem code \texttt{catalogue entry}, title ``Pure Resolutions of Paracanonical Level Curves.'' The supplied record gives proposed year 2012 and names no proposer.
\item Supplied report citation: \emph{Classical Algebraic Geometry (2013)}, Oberwolfach Reports, pp.~1845--1893, DOI \texttt{10.4171/owr/2012/30}, \url{https://doi.org/10.4171/owr/2012/30}. The background describes Oberwolfach Reports as documenting mathematical workshops and their open-problem discussions. The citation's 2013 label differs from its DOI year and the recorded proposed year 2012; this has not been reconciled.
\item Supplied literature reference: Farkas, \emph{Syzygies of torsion bundles}, \url{https://www2.mathematik.hu-berlin.de/~farkas/nivel3.pdf}. No further bibliographic details are supplied.
\item Supplied literature reference labeled ``original OWR report'': \url{https://doi.org/10.4171/owr/2012/30}.
\item Literature assessment dated exactly \texttt{2026-08-21}: the supplied curation labels the problem open and reports verification in low genus and special levels, with exceptional genera identified, while reporting that no general proof covering the stated conjectural pattern was found. This is supplied dated curation derived from the OWR material, not an independently verified literature assessment; the status and cited claims require independent checking.
\item The supplied statement status is \texttt{corrected\_verified}; its verification note reports comparison with the official Oberwolfach report and normalization of notation, context, and question form without changing the source-level claim. This is a supplied verification claim, not verification performed for this draft.
\end{itemize}
\end{document}
